\documentclass[12pt,a4paper]{article}

\usepackage[utf8]{inputenc}
\usepackage[T1]{fontenc}
\usepackage[french]{babel}

\usepackage{styles/couleurs}
\usepackage{styles/mise_en_page}
\usepackage{styles/boites}
\usepackage{styles/maths}
\usepackage{styles/figures}

\usepackage{amsmath}
\usepackage{amssymb}
\usepackage{enumitem}
\usepackage{array}
\usepackage{tabularx}
\usepackage{graphicx}
\usepackage{tikz}
\usepackage{xparse}
\usetikzlibrary{calc,angles,quotes}

\setlist[itemize]{leftmargin=1.4em,itemsep=0.15em,topsep=0.2em}
\setlist[enumerate]{leftmargin=*,itemsep=0.22em,topsep=0.22em}

\newcommand{\titresequence}[1]{%
\begin{center}
  {\sffamily\LARGE\bfseries\color{TitreC}Séquence 2 -- Théorème de Pythagore}\\[0.2em]
  {\sffamily\large\bfseries\color{SavoirC}#1}
\end{center}}

\newcommand{\lignepointillee}{\noindent\dotfill\par\vspace{0.45em}}
\newcommand{\lignesreponse}[1]{%
\begin{tcolorbox}[
  enhanced,colback=white,colframe=NeutreC!40,boxrule=0.5pt,arc=1.5mm,
  left=3mm,right=3mm,top=2mm,bottom=1mm,before skip=2pt,after skip=2pt]
  \foreach \i in {1,...,#1}{\lignepointillee}
\end{tcolorbox}}

\begin{document}
\titresequence{Activité 0 -- Aire, périmètre et côté d'un carré}

Cette activité est courte. Elle sert seulement à réactiver les idées indispensables au puzzle.

\begin{enumerate}
  \item Dessine deux figures formées de carreaux qui ont la même aire mais des périmètres différents.

  \item Dessine deux figures formées de carreaux qui ont le même périmètre mais des aires différentes.

  \item Complète :
  \begin{enumerate}[label=\alph*)]
    \item Le périmètre d'un carré vaut $36$ cm. Son côté mesure \dotfill
    \item L'aire d'un carré vaut $36$ cm$^2$. Son côté mesure \dotfill
    \item L'aire d'un carré vaut $49$ cm$^2$. Son côté mesure \dotfill
  \end{enumerate}
\end{enumerate}

\begin{tcolorbox}[colback=SavoirC!5,colframe=SavoirC!60,boxrule=0.6pt,arc=2mm]

\textbf{Bilan.}
\lignesreponse{5}
\end{tcolorbox}

\hspace{0.5cm}

\titresequence{Activité 0 -- Aire, périmètre et côté d'un carré}

Cette activité est courte. Elle sert seulement à réactiver les idées indispensables au puzzle.

\begin{enumerate}
  \item Dessine deux figures formées de carreaux qui ont la même aire mais des périmètres différents.

  \item Dessine deux figures formées de carreaux qui ont le même périmètre mais des aires différentes.

  \item Complète :
  \begin{enumerate}[label=\alph*)]
    \item Le périmètre d'un carré vaut $36$ cm. Son côté mesure \dotfill
    \item L'aire d'un carré vaut $36$ cm$^2$. Son côté mesure \dotfill
    \item L'aire d'un carré vaut $49$ cm$^2$. Son côté mesure \dotfill
  \end{enumerate}
\end{enumerate}

\begin{tcolorbox}[colback=SavoirC!5,colframe=SavoirC!60,boxrule=0.6pt,arc=2mm]

\textbf{Bilan.}
\lignesreponse{5}
\end{tcolorbox}

\newpage
\textbf{Bilan.}
Le périmètre mesure le contour d'une figure. L'aire mesure la surface intérieure.
Si un carré a pour côté $c$, son aire vaut $c^2$.
Chercher le côté d'un carré à partir de son aire revient à chercher le nombre dont le carré est cette aire.
\end{document}
